\documentclass[12pt]{article}
\usepackage[utf8]{inputenc}
\usepackage[T1]{fontenc}
\usepackage{hyperref}
\title{Mituku Language Resources for K-12 Curriculum Development}
\author{Compiled by Claude Sonnet 5.5 (claude-sonnet-5-5) for Andr\'as Kornai}
\date{}
\begin{document}
\maketitle

\section{Introduction}
Mituku (ISO 639-3 zmq, Glottocode mitu1240) is a Bantu language of the Democratic Republic of the Congo. It is also called Kinya-Mituku or Metoko. It is tonal, with four tones. Sources place it in the Ubundu area of the former Orientale Province, west of the Lualaba River. Speaker estimates range from 51,000 to 117,000. Search found no dictionary, primer or school materials.

\section{Primary Educational Resources}
None were found. The language is not known to be taught in schools.

\section{Linguistic and Reference Materials}
\begin{itemize}
\item Wikipedia, Mituku language: \url{https://en.wikipedia.org/wiki/Mituku_language}
\item ASJP, Mituku word list: \url{https://asjp.clld.org/languages/MITUKU}
\end{itemize}

\section{Religious and Community Resources}
\begin{itemize}
\item Joshua Project, Kinya-Mituku in the DR Congo (lists scripture translation as started): \url{https://www.joshuaproject.net/index.php/people_groups/12703}
\end{itemize}

\section{Digital and Audio Resources}
None were found.

\section{Curriculum Development Context}
Mituku has a sizeable speaker community, but almost nothing is documented. Curriculum work would need new description of the orthography and lexicon, and community involvement. The scripture translation that Joshua Project reports as started could become a source of text if it is published.

\end{document}